%===========================================================
% 05_geometry.tex - 几何变换
%===========================================================

\section{几何变换}

%-----------------------------------------------------------
% Frame 1: 仿射变换基础
%-----------------------------------------------------------
\begin{frame}{仿射变换基础}
    \begin{columns}
        \column{0.5\textwidth}
        \begin{block}{定义}
            \begin{itemize}
                \item 保持\textbf{平行性}的变换
                \item 直线变换后仍是直线
                \item 平行线保持平行
            \end{itemize}
        \end{block}

        \begin{block}{包含操作}
            \begin{itemize}
                \item 平移 (Translation)
                \item 旋转 (Rotation)
                \item 缩放 (Scale)
                \item 倾斜 (Shear)
            \end{itemize}
        \end{block}

        \column{0.5\textwidth}
        \begin{block}{变换矩阵}
            $$M = \begin{bmatrix}
            a_{11} & a_{12} & b_x \\
            a_{21} & a_{22} & b_y
            \end{bmatrix}$$

            \begin{itemize}
                \item 2×3 矩阵
                \item 6 个参数确定
                \item 用 \code{cv2.warpAffine()}
            \end{itemize}
        \end{block}
    \end{columns}

    \begin{alertblock}{重要特性}
        仿射变换只需要\textbf{3个点}确定（6个方程）
    \end{alertblock}
\end{frame}

%-----------------------------------------------------------
% Frame 2: 仿射变换代码
%-----------------------------------------------------------
\begin{frame}[fragile]{仿射变换代码}
    \begin{columns}
        \column{0.5\textwidth}
        \begin{lstlisting}[basicstyle=\ttfamily\scriptsize]
import cv2
import numpy as np

# 读取图像
img = cv2.imread('exam.jpg')
h, w = img.shape[:2]

# 1. 平移变换
M_trans = np.float32([
    [1, 0, 50],    # x 方向平移 50
    [0, 1, 30]     # y 方向平移 30
])
translated = cv2.warpAffine(img, M_trans, (w, h))

# 2. 旋转变换
center = (w // 2, h // 2)
angle = 45
scale = 1.0
M_rot = cv2.getRotationMatrix2D(center, angle, scale)
rotated = cv2.warpAffine(img, M_rot, (w, h))

# 3. 组合变换
M = np.float32([
    [1.2, 0.2, 0],
    [0, 1.1, 0]
])
scaled = cv2.warpAffine(img, M, (w, h))
        \end{lstlisting}

        \column{0.5\textwidth}
        \begin{block}{常用函数}
            \begin{itemize}
                \item \code{cv2.warpAffine()}：应用变换
                \item \code{cv2.getRotationMatrix2D()}：获取旋转矩阵
                \item \code{cv2.getAffineTransform()}：从3点计算矩阵
            \end{itemize}
        \end{block}

        \begin{alertblock}{TODO 任务}
            实现一个函数：\\
            输入图像 + 角度，返回旋转后的图像
        \end{alertblock}
    \end{columns}
\end{frame}

%-----------------------------------------------------------
% Frame 3: 透视变换
%-----------------------------------------------------------
\begin{frame}[fragile]{透视变换}
    \textbf{原理：} 4点对应变换，可以矫正倾斜的矩形

    \begin{columns}
        \column{0.5\textwidth}
        \begin{block}{变换矩阵}
            \begin{itemize}
                \item 需要 \highlight{4个点}确定
                \item 3×3 变换矩阵
                \item 用 \code{cv2.getPerspectiveTransform()}
                \item 用 \code{cv2.warpPerspective()}
            \end{itemize}
        \end{block}

        \column{0.5\textwidth}
        \begin{lstlisting}[basicstyle=\ttfamily\scriptsize]
import cv2
import numpy as np

# 原图上的4个点（顺时针）
pts1 = np.float32([
    [100, 50],     # 左上
    [400, 50],     # 右上
    [400, 300],    # 右下
    [100, 300]     # 左下
])

# 希望变换到的位置
pts2 = np.float32([
    [0, 0],
    [300, 0],
    [300, 250],
    [0, 250]
])

# 计算透视变换矩阵
M = cv2.getPerspectiveTransform(pts1, pts2)

# 应用变换
h, w = 300, 300
result = cv2.warpPerspective(img, M, (w, h))

# 显示
cv2.imshow('Perspective', result)
cv2.waitKey(0)
        \end{lstlisting}
    \end{columns}

    \begin{alertblock}{应用场景}
        试卷倾斜矫正、文档扫描、透视畸变校正
    \end{alertblock}
\end{frame}

%-----------------------------------------------------------
% Frame 4: 自动矫正实战
%-----------------------------------------------------------
\begin{frame}[fragile]{自动矫正实战：试卷倾斜矫正}
	\begin{alertblock}{GUI演示程序：\code{gui\_04\_geometry.py}}
		运行 \code{python gui\_launcher.py} 选择“几何变换演示”，可交互调节旋转角度、缩放和平移参数
	\end{alertblock}

    \begin{lstlisting}[basicstyle=\ttfamily\scriptsize]
import cv2
import numpy as np

def correct_perspective(image_path):
    """
    试卷倾斜矫正完整流程
    """
    # Step 1: 读取图像
    img = cv2.imread(image_path)
    gray = cv2.cvtColor(img, cv2.COLOR_BGR2GRAY)

    # Step 2: 边缘检测
    edges = cv2.Canny(gray, 50, 150, apertureSize=3)

    # Step 3: 霍夫变换检测直线
    lines = cv2.HoughLinesP(edges, 1, np.pi/180, 100,
                            minLineLength=100, maxLineGap=10)

    # Step 4: 计算角度
    angles = []
    for line in lines:
        x1, y1, x2, y2 = line[0]
        angle = np.arctan2(y2-y1, x2-x1) * 180.0 / np.pi
        angles.append(angle)

    # Step 5: 取中位数角度
    median_angle = np.median(angles)

    # Step 6: 旋转矫正
    h, w = img.shape[:2]
    center = (w // 2, h // 2)
    M = cv2.getRotationMatrix2D(center, median_angle, 1.0)
    corrected = cv2.warpAffine(img, M, (w, h),
                               flags=cv2.INTER_CUBIC)

    return corrected

# 测试
result = correct_perspective('exam.jpg')
cv2.imshow('Corrected', result)
cv2.waitKey(0)
    \end{lstlisting}

    \begin{alertblock}{关键步骤}
        \begin{enumerate}
            \item 边缘检测找轮廓
            \item 霍夫变换检测直线
            \item 计算倾斜角度
            \item 旋转变换矫正
        \end{enumerate}
    \end{alertblock}
\end{frame}

%-----------------------------------------------------------
% Frame 5: 挑战任务
%-----------------------------------------------------------
\begin{frame}[fragile]{挑战任务：智能裁剪答题卡}
    \begin{columns}
        \column{0.5\textwidth}
        \begin{block}{任务描述}
            实现一个函数，输入试卷图片，自动找到答题卡区域并裁剪

            \begin{enumerate}
                \item 检测最大轮廓（可能是答题卡）
                \item 获取4个角点
                \item 透视变换裁剪
                \item 返回裁剪后的图像
            \end{enumerate}
        \end{block}

        \column{0.5\textwidth}
        \begin{lstlisting}[basicstyle=\ttfamily\scriptsize]
import cv2
import numpy as np

def crop_answer_card(image_path):
    """智能裁剪答题卡"""
    img = cv2.imread(image_path)
    gray = cv2.cvtColor(img, cv2.COLOR_BGR2GRAY)

    # TODO Step 1: 二值化
    # ...

    # TODO Step 2: 找轮廓
    # contours, _ = cv2.findContours(...)

    # TODO Step 3: 找最大四边形
    # hint: cv2.approxPolyDP

    # TODO Step 4: 透视变换裁剪
    # ...

    return cropped

# 测试
result = crop_answer_card('exam.jpg')
cv2.imshow('Cropped', result)
cv2.waitKey(0)
        \end{lstlisting}
    \end{columns}

    \begin{alertblock}{提示}
        \begin{itemize}
            \item 先用 Canny 边缘检测
            \item 用 \code{cv2.findContours} 找轮廓
            \item 用 \code{cv2.approxPolyDP} 近似多边形
            \item 保留有4个顶点的轮廓
        \end{itemize}
    \end{alertblock}
\end{frame}
